\section{Experimental Method}
\label{sec:method}

\subsection{System Configuration}
\begin{chinese}
\subsection*{系统配置}
\end{chinese}

All experiments are conducted through the Pigs HTTP API. The system is configured with:
\begin{chinese}
所有实验通过 Pigs HTTP API 进行。系统配置为：
\end{chinese}
\begin{itemize}[nosep]
    \item Default orchestration limits: \texttt{max\_post\_iterations=3}, \texttt{max\_pre\_replans=2}
    \begin{chinese}
    默认编排限制：\texttt{max\_post\_iterations=3}，\texttt{max\_pre\_replans=2}
    \end{chinese}
    \item Continuation TTL: 30 minutes, capacity: 256
    \begin{chinese}
    Continuation TTL：30 分钟，容量：256
    \end{chinese}
    \item Language: Chinese (zh) for phase prompts
    \begin{chinese}
    语言：中文（zh）相位提示词
    \end{chinese}
    \item Thinking effort: xhigh (OpenAI) / max (Anthropic)
    \begin{chinese}
    思考强度：xhigh（OpenAI）/ max（Anthropic）
    \end{chinese}
\end{itemize}

\subsection{Ablation Conditions}
\begin{chinese}
\subsection*{消融条件}
\end{chinese}

For Experiment 1, we define four conditions:
\begin{chinese}
实验 1 定义四个条件：
\end{chinese}
\begin{itemize}
    \item \textbf{Condition A (Baseline)}: Direct API call without orchestration (non-\texttt{-pig} model, passthrough).
    \begin{chinese}
    \textbf{条件 A（基线）}：直接 API 调用，无编排（非 \texttt{-pig} 模型，透传）。
    \end{chinese}
    \item \textbf{Condition B (Full)}: Complete Pre$\rightarrow$Executor$\rightarrow$Post orchestration (\texttt{-pig} model).
    \begin{chinese}
    \textbf{条件 B（完整）}：完整 Pre$\rightarrow$Executor$\rightarrow$Post 编排（\texttt{-pig} 模型）。
    \end{chinese}
    \item \textbf{Condition C (Pre-only)}: Only the Pre phase (planning without execution verification).
    \begin{chinese}
    \textbf{条件 C（仅 Pre）}：只有 Pre 阶段（规划但不验证执行）。
    \end{chinese}
    \item \textbf{Condition D (Executor+Post)}: Skip Pre, directly execute and review.
    \begin{chinese}
    \textbf{条件 D（Executor+Post）}：跳过 Pre，直接执行和评审。
    \end{chinese}
\end{itemize}

\subsection{Metrics}
\begin{chinese}
\subsection*{评估指标}
\end{chinese}

\begin{itemize}[nosep]
    \item \textbf{Task completion rate}: percentage of tasks successfully completed (benchmark-specific definition).
    \begin{chinese}
    \textbf{任务完成率}：成功完成的任务百分比（基准特定定义）。
    \end{chinese}
    \item \textbf{Average turns}: mean number of LLM API calls per task.
    \begin{chinese}
    \textbf{平均轮次}：每个任务的 LLM API 调用次数均值。
    \end{chinese}
    \item \textbf{Token consumption}: total input + output tokens per task.
    \begin{chinese}
    \textbf{Token 消耗}：每个任务的总输入 + 输出 token 数。
    \end{chinese}
    \item \textbf{Latency}: wall-clock time from request to final response.
    \begin{chinese}
    \textbf{延迟}：从请求到最终响应的实际时间。
    \end{chinese}
    \item \textbf{Error recovery rate}: percentage of tasks where PIGFAIL triggered a successful replan.
    \begin{chinese}
    \textbf{错误恢复率}：PIGFAIL 触发成功重规划的任务百分比。
    \end{chinese}
\end{itemize}

\subsection{Benchmarks}
\begin{chinese}
\subsection*{基准测试}
\end{chinese}

\begin{center}
\begin{tabular}{lllr}
\toprule
Benchmark & Domain & Tasks & Public \\
\midrule
SWE-bench Lite~\citep{jimenez2024swebench} & Software engineering & 300 & \checkmark \\
GAIA~\citep{mialon2023gaia} & General AI assistant & 300 & \checkmark \\
AgentBoard~\citep{ma2024agentboard} & Multi-turn agent & 9 tasks & \checkmark \\
\bottomrule
\end{tabular}
\end{center}
